% Generated from the matching Typst scaffold by
% scripts/migrate_typst_report_to_latex.py (prose drafted 2026-06-25).

% ── 4.3  Memory Architecture ─────────────────────────────────────────────
% PURPOSE: Describe the storage substrate every policy shares.
% ─────────────────────────────────────────────────────────────────────────

\section{Memory Architecture}

All six conditions share one storage substrate; the policies differ only in how they prune
it. Each completed task can write a \texttt{MemoryRecord} into a \texttt{MemoryStore} backed
by a FAISS index, and every condition reads and writes the same store through the same
interface (Figure~\ref{fig:architecture}). What the policies later act on is the
\emph{footprint} of this store --- the number and size of records carried forward.

The single most consequential design choice is what gets embedded. Rather than embedding the
raw trajectory --- which the embedder would silently truncate at its token cap, turning a
long agent transcript into a meaningless prefix --- each record embeds a compact,
outcome-bearing triple: the issue text, the final error, and the final diff. This payload is
capped under 7500 tokens and the bound is enforced at write time
(\texttt{store.py::\_verify\_embedding\_size}, Invariant~\#4), so a record can never silently
exceed the embedder's window. The compactness is deliberate: it keeps the embedding faithful
to the task's outcome and keeps the per-record footprint small.

Each record also carries a \emph{content}-based type drawn from a five-type taxonomy ---
\texttt{architectural}, \texttt{api\_change}, \texttt{bug\_fix}, \texttt{test\_update},
\texttt{config} (Invariant~\#7). The type describes what the record is \emph{about}, not
whether the task succeeded; \texttt{outcome} is a separate, orthogonal axis and is never
collapsed into the type. A dedicated classifier assigns the type at write time (its
structured-output handling under the substituted model is disclosed as deviation D4/D6b in
Section~\ref{sec:deviations}).

The write path is part of the shared per-task loop (Figure~\ref{fig:task-lifecycle}): after
the agent attempts the task, a reflection step produces one record, the classifier labels
it, and the maintenance step --- the only policy-specific step --- decides what the store
keeps. Type-Aware Decay reads the type at this step; the other policies ignore it. The cap of
ten (A1) bounds the store for the four interior policies, and that bounded footprint is the
quantity the cost analysis in Section~\ref{sec:metrics} tracks.
